\documentclass[11pt]{article}
\usepackage[a4paper,margin=2.2cm]{geometry}
\usepackage{amsmath,amssymb,booktabs,longtable}
\usepackage[T1]{fontenc}
\usepackage{hyperref}
% generated by qbscreen/make_numbers_tex.py — do not edit by hand
\newcommand{\tContactMin}{109}
\newcommand{\tContactMax}{290}
\newcommand{\tContactGeo}{160}
\newcommand{\betaMin}{2.5}
\newcommand{\betaMax}{3.4}
\newcommand{\nScanPoints}{33}
\newcommand{\lamFlHalf}{0.20}
\newcommand{\lamMethyl}{0.63}
\newcommand{\lamBenzyl}{0.48}
\newcommand{\lamAla}{0.92}
\newcommand{\aNfiveDFT}{441}
\newcommand{\aNtenDFT}{165}
\newcommand{\partnerMAO}{257}
\newcommand{\partnerDAO}{58}
\newcommand{\DeltaMaxMin}{0.53}
\newcommand{\DeltaMaxMax}{0.93}
\newcommand{\tauCompMax}{\ensuremath{3.9\times10^{-11}}}
\newcommand{\tauCompMin}{\ensuremath{1.1\times10^{-13}}}
\newcommand{\JCompMin}{\ensuremath{2.9\times10^{5}}}
\newcommand{\JCompMax}{\ensuremath{5.6\times10^{7}}}
\newcommand{\JtauInverted}{57}
\newcommand{\JtauActless}{0.18}
\newcommand{\DMAO}{-1817}
\newcommand{\DDAO}{-1217}
\newcommand{\capMAO}{\ensuremath{5.0\times10^{-5}}}
\newcommand{\capDAO}{\ensuremath{9.3\times10^{-6}}}
\newcommand{\tAnchor}{22}
\newcommand{\tAnchorWb}{4}
\newcommand{\partnerCapMin}{\ensuremath{3.6\times10^{-8}}}
\newcommand{\partnerCapMax}{\ensuremath{8.8\times10^{-5}}}
\newcommand{\spinOneRatio}{2.0}
\newcommand{\DMAOspin}{-914}
\newcommand{\DDAOspin}{-674}
\newcommand{\capMAOspin}{\ensuremath{4.3\times10^{-4}}}
\newcommand{\capDAOspin}{\ensuremath{4.1\times10^{-4}}}
\newcommand{\DCRY}{-11.4}
\newcommand{\cryNinetyMin}{0.21}
\newcommand{\cryNinetyMax}{1.74}
\newcommand{\cryPowderMin}{0.18}
\newcommand{\cryPowderMax}{1.42}

\renewcommand{\thetable}{S\arabic{table}}
\renewcommand{\thesection}{S\arabic{section}}

\title{Supporting Information:\\ A thermally formed flavin radical pair that forms
fast enough to matter cannot sense the geomagnetic field}
\author{Hikaru Wakaura and Taiki Tanimae}
\date{}

\begin{document}
\maketitle

\section{Software and environment}
Spin dynamics: \texttt{qbscreen} (this work), exact Liouville-space solution with
NumPy. Electronic structure: PySCF~2.12.1 (B3LYP, def2-SVP and def2-TZVP, RI-J
density fitting); geometries: xtb~6.7.1 (GFN2-xTB, L-BFGS optimizer, tight
convergence). Every number in the main text and in the tables below is produced by
the scripts \texttt{qbscreen/calibration\_run.py},
\texttt{qbscreen/calibration\_extras.py}, \texttt{qbscreen/make\_numbers\_tex.py}
and \texttt{qbscreen/make\_si\_tables.py} from the stored results in
\texttt{calibration\_results/}.

\section{Electronic coupling scan}
Methylamine was placed on the si face of lumiflavin with its nitrogen on the ring
normal through N5 at distance $r$ and its lone pair directed at N5. Twist rotates
the amine about that axis; tilt tips the axis off the normal. Geometries whose
closest intermolecular contact is below 2.4~\AA\ are steric clashes of the rigid
placement and were excluded from the fits (Table~\ref{tab:scan}).

\begin{table}[h]\centering\small
\caption{Fragment-projection coupling $t$ between the flavin LUMO and the amine HOMO
(B3LYP/def2-SVP).}\label{tab:scan}
\begin{tabular}{rrrrrl}
\toprule
$r$ (\AA) & twist ($^\circ$) & tilt ($^\circ$) & $t$ (meV) & closest contact (\AA) & used \\
\midrule
3.00 & 0 & -15 & 284.9 & 2.84 & yes \\
3.00 & 0 & 0 & 615.4 & 2.50 & yes \\
3.00 & 0 & 15 & 882.7 & 1.99 & no (clash) \\
3.00 & 60 & 0 & 248.1 & 2.51 & yes \\
3.25 & 0 & -15 & 177.3 & 3.09 & yes \\
3.25 & 0 & 0 & 424.8 & 2.75 & yes \\
3.25 & 0 & 15 & 623.0 & 2.18 & no (clash) \\
3.25 & 60 & 0 & 179.3 & 2.76 & yes \\
3.50 & 0 & -15 & 109.1 & 3.34 & yes \\
3.50 & 0 & 0 & 290.4 & 3.00 & yes \\
3.50 & 0 & 15 & 428.6 & 2.38 & no (clash) \\
3.50 & 60 & 0 & 129.6 & 3.01 & yes \\
3.75 & 0 & -15 & 66.7 & 3.56 & yes \\
3.75 & 0 & 0 & 197.2 & 3.25 & yes \\
3.75 & 0 & 15 & 290.5 & 2.58 & yes \\
3.75 & 60 & 0 & 93.6 & 3.23 & yes \\
4.00 & 0 & -15 & 40.9 & 3.79 & yes \\
4.00 & 0 & 0 & 133.5 & 3.50 & yes \\
4.00 & 0 & 15 & 195.2 & 2.80 & yes \\
4.00 & 60 & 0 & 67.7 & 3.41 & yes \\
4.50 & 0 & -15 & 16.6 & 4.26 & yes \\
4.50 & 0 & 0 & 62.6 & 4.00 & yes \\
4.50 & 0 & 15 & 88.2 & 3.24 & yes \\
4.50 & 60 & 0 & 36.8 & 3.81 & yes \\
5.00 & 0 & -15 & 8.0 & 4.73 & yes \\
5.00 & 0 & 0 & 31.0 & 4.50 & yes \\
5.00 & 0 & 15 & 42.7 & 3.70 & yes \\
5.00 & 60 & 0 & 21.2 & 4.23 & yes \\
5.50 & 0 & -15 & 4.0 & 5.21 & yes \\
5.50 & 0 & 0 & 15.1 & 5.00 & yes \\
5.50 & 0 & 15 & 21.9 & 4.17 & yes \\
5.50 & 60 & 0 & 11.5 & 4.66 & yes \\
6.00 & 0 & -15 & 1.8 & 5.69 & yes \\
6.00 & 0 & 0 & 6.5 & 5.50 & yes \\
6.00 & 0 & 15 & 10.5 & 4.64 & yes \\
6.00 & 60 & 0 & 5.3 & 5.10 & yes \\
\bottomrule
\end{tabular}

\end{table}

\section{Inner-sphere reorganization energy}
Four-point energies (hartree) of each redox couple, state $s$ at the optimized
geometry $G$ of state $s'$, and the resulting half contributions
$\lambda_i/2 = \tfrac12\{[E_0(G_1)-E_0(G_0)] + [E_1(G_0)-E_1(G_1)]\}$
(Table~\ref{tab:lambda}). State 0 is the closed-shell species, state 1 the radical
ion.

\begin{table}[h]\centering\small
\caption{Four-point energies (B3LYP/def2-SVP at GFN2-xTB geometries).}\label{tab:lambda}
\begin{tabular}{lrrrrr}
\toprule
species & $E_0(G_0)$ & $E_0(G_1)$ & $E_1(G_1)$ & $E_1(G_0)$ & $\lambda_i/2$ (eV) \\
\midrule
lumiflavin & -871.480317 & -871.473953 & -871.539794 & -871.531658 & 0.197 \\
methylamine & -95.780457 & -95.771484 & -95.461713 & -95.439105 & 0.430 \\
benzylamine & -326.673644 & -326.663153 & -326.374740 & -326.364444 & 0.283 \\
alanine & -323.509936 & -323.476155 & -323.178642 & -323.158966 & 0.727 \\
\bottomrule
\end{tabular}

\end{table}

\section{Hyperfine couplings}
Isotropic Fermi-contact couplings from the UKS spin density at each nucleus
(B3LYP/def2-TZVP), six largest per radical (Table~\ref{tab:hfc}). The flavin
anion N5 and N10 values, \aNfiveDFT\ and \aNtenDFT~$\mu$T, compare with 523 and
189~$\mu$T of Lee et al.; the spin model uses the latter. The large aminium
couplings come from $\alpha$-C--H protons hyperconjugated with the
nitrogen-centred orbital and depend on conformation.

\begin{table}[h]\centering\small
\caption{DFT isotropic hyperfine couplings.}\label{tab:hfc}
\begin{tabular}{llrr}
\toprule
radical & nucleus & $a_{\rm iso}$ (MHz) & $a_{\rm iso}$ (mT) \\
\midrule
flavin$^{\bullet-}$ & H29 & +17.9 & +0.639 \\
flavin$^{\bullet-}$ & H30 & +17.8 & +0.636 \\
flavin$^{\bullet-}$ & N4 & +12.4 & +0.441 \\
flavin$^{\bullet-}$ & H25 & +11.3 & +0.403 \\
flavin$^{\bullet-}$ & H26 & +11.3 & +0.402 \\
flavin$^{\bullet-}$ & H22 & -10.5 & -0.374 \\
benzylaminium$^{\bullet+}$ & H10 & +256.9 & +9.167 \\
benzylaminium$^{\bullet+}$ & H11 & +196.2 & +7.001 \\
benzylaminium$^{\bullet+}$ & N0 & +24.8 & +0.885 \\
benzylaminium$^{\bullet+}$ & H8 & -23.0 & -0.821 \\
benzylaminium$^{\bullet+}$ & H9 & -20.7 & -0.737 \\
benzylaminium$^{\bullet+}$ & H14 & -12.3 & -0.439 \\
methylaminium$^{\bullet+}$ & H2 & +363.1 & +12.955 \\
methylaminium$^{\bullet+}$ & H6 & +82.7 & +2.951 \\
methylaminium$^{\bullet+}$ & H5 & +81.9 & +2.923 \\
methylaminium$^{\bullet+}$ & H3 & -59.3 & -2.116 \\
methylaminium$^{\bullet+}$ & H4 & -59.3 & -2.115 \\
methylaminium$^{\bullet+}$ & N1 & +30.5 & +1.090 \\
alaninium$^{\bullet+}$ & H9 & +58.3 & +2.080 \\
alaninium$^{\bullet+}$ & H6 & +30.8 & +1.098 \\
alaninium$^{\bullet+}$ & N2 & +18.6 & +0.665 \\
alaninium$^{\bullet+}$ & H11 & -17.9 & -0.639 \\
alaninium$^{\bullet+}$ & H10 & -17.7 & -0.632 \\
alaninium$^{\bullet+}$ & H7 & +14.1 & +0.502 \\
\bottomrule
\end{tabular}

\end{table}

\section{Catalytic-competence bound: full grid}
\begin{longtable}{c}
\caption{Upper bound on $\Delta$ for thermal formation at $k_{\rm cat}$, with the
lifetime and exchange coupling at the bound.}\label{tab:competence}\\
\begin{tabular}{rrrrrr}
\toprule
$t$ (meV) & $\lambda$ (eV) & $k_{\rm cat}$ (s$^{-1}$) & $\Delta_{\max}$ (eV) & $\tau$ (s) & $|J|$ (MHz) \\
\midrule
22 & 0.7 & 1 & 0.80 & $\ensuremath{2.3\times10^{-13}}$ & $\ensuremath{2.9\times10^{5}}$ \\
22 & 0.7 & 10 & 0.74 & $\ensuremath{2.1\times10^{-13}}$ & $\ensuremath{3.2\times10^{5}}$ \\
22 & 0.7 & 100 & 0.68 & $\ensuremath{2.1\times10^{-13}}$ & $\ensuremath{3.5\times10^{5}}$ \\
22 & 1.0 & 1 & 0.78 & $\ensuremath{3.5\times10^{-13}}$ & $\ensuremath{3.0\times10^{5}}$ \\
22 & 1.0 & 10 & 0.71 & $\ensuremath{4.9\times10^{-13}}$ & $\ensuremath{3.3\times10^{5}}$ \\
22 & 1.0 & 100 & 0.64 & $\ensuremath{7.7\times10^{-13}}$ & $\ensuremath{3.7\times10^{5}}$ \\
22 & 1.4 & 1 & 0.70 & $\ensuremath{6.4\times10^{-12}}$ & $\ensuremath{3.3\times10^{5}}$ \\
22 & 1.4 & 10 & 0.62 & $\ensuremath{1.5\times10^{-11}}$ & $\ensuremath{3.8\times10^{5}}$ \\
22 & 1.4 & 100 & 0.53 & $\ensuremath{3.9\times10^{-11}}$ & $\ensuremath{4.4\times10^{5}}$ \\
109 & 0.7 & 1 & 0.87 & $\ensuremath{1.6\times10^{-13}}$ & $\ensuremath{6.6\times10^{6}}$ \\
109 & 0.7 & 10 & 0.82 & $\ensuremath{1.3\times10^{-13}}$ & $\ensuremath{7.0\times10^{6}}$ \\
109 & 0.7 & 100 & 0.76 & $\ensuremath{1.1\times10^{-13}}$ & $\ensuremath{7.6\times10^{6}}$ \\
109 & 1.0 & 1 & 0.88 & $\ensuremath{1.2\times10^{-13}}$ & $\ensuremath{6.6\times10^{6}}$ \\
109 & 1.0 & 10 & 0.81 & $\ensuremath{1.5\times10^{-13}}$ & $\ensuremath{7.1\times10^{6}}$ \\
109 & 1.0 & 100 & 0.74 & $\ensuremath{2.0\times10^{-13}}$ & $\ensuremath{7.8\times10^{6}}$ \\
109 & 1.4 & 1 & 0.81 & $\ensuremath{1.1\times10^{-12}}$ & $\ensuremath{7.1\times10^{6}}$ \\
109 & 1.4 & 10 & 0.73 & $\ensuremath{2.1\times10^{-12}}$ & $\ensuremath{7.8\times10^{6}}$ \\
109 & 1.4 & 100 & 0.65 & $\ensuremath{4.5\times10^{-12}}$ & $\ensuremath{8.8\times10^{6}}$ \\
160 & 0.7 & 1 & 0.89 & $\ensuremath{1.7\times10^{-13}}$ & $\ensuremath{1.4\times10^{7}}$ \\
160 & 0.7 & 10 & 0.84 & $\ensuremath{1.3\times10^{-13}}$ & $\ensuremath{1.5\times10^{7}}$ \\
160 & 0.7 & 100 & 0.78 & $\ensuremath{1.1\times10^{-13}}$ & $\ensuremath{1.6\times10^{7}}$ \\
160 & 1.0 & 1 & 0.90 & $\ensuremath{1.1\times10^{-13}}$ & $\ensuremath{1.4\times10^{7}}$ \\
160 & 1.0 & 10 & 0.83 & $\ensuremath{1.3\times10^{-13}}$ & $\ensuremath{1.5\times10^{7}}$ \\
160 & 1.0 & 100 & 0.76 & $\ensuremath{1.7\times10^{-13}}$ & $\ensuremath{1.6\times10^{7}}$ \\
160 & 1.4 & 1 & 0.84 & $\ensuremath{8.4\times10^{-13}}$ & $\ensuremath{1.5\times10^{7}}$ \\
160 & 1.4 & 10 & 0.76 & $\ensuremath{1.6\times10^{-12}}$ & $\ensuremath{1.6\times10^{7}}$ \\
160 & 1.4 & 100 & 0.68 & $\ensuremath{3.3\times10^{-12}}$ & $\ensuremath{1.8\times10^{7}}$ \\
290 & 0.7 & 1 & 0.92 & $\ensuremath{1.9\times10^{-13}}$ & $\ensuremath{4.4\times10^{7}}$ \\
290 & 0.7 & 10 & 0.86 & $\ensuremath{1.4\times10^{-13}}$ & $\ensuremath{4.7\times10^{7}}$ \\
290 & 0.7 & 100 & 0.81 & $\ensuremath{1.2\times10^{-13}}$ & $\ensuremath{5.0\times10^{7}}$ \\
290 & 1.0 & 1 & 0.93 & $\ensuremath{1.1\times10^{-13}}$ & $\ensuremath{4.4\times10^{7}}$ \\
290 & 1.0 & 10 & 0.87 & $\ensuremath{1.2\times10^{-13}}$ & $\ensuremath{4.7\times10^{7}}$ \\
290 & 1.0 & 100 & 0.80 & $\ensuremath{1.5\times10^{-13}}$ & $\ensuremath{5.1\times10^{7}}$ \\
290 & 1.4 & 1 & 0.88 & $\ensuremath{6.2\times10^{-13}}$ & $\ensuremath{4.6\times10^{7}}$ \\
290 & 1.4 & 10 & 0.80 & $\ensuremath{1.1\times10^{-12}}$ & $\ensuremath{5.1\times10^{7}}$ \\
290 & 1.4 & 100 & 0.72 & $\ensuremath{2.2\times10^{-12}}$ & $\ensuremath{5.6\times10^{7}}$ \\
\bottomrule
\end{tabular}

\end{longtable}

\section{Magnetic field effect along the coupling-locked locus}
For the MAO contact pair ($r = 0.35$~nm, point-dipole $D = \DMAO$~MHz), the
exchange coupling and lifetime follow from $t$ through equations~2--4 of the main
text, with the adiabatic interpolation. MFE in \% at 50~$\mu$T, coherent and with
$T_2^e = 1~\mu$s, field at 0 and 90$^\circ$ to the interradical axis.

\begin{table}[h]\centering\footnotesize
\caption{Inverted region, $\Delta = 1.78$, $\lambda = 0.97$~eV.}\label{tab:locinv}
\begin{tabular}{rrrrrrr}
\toprule
$t$ (meV) & $|J|$ (MHz) & $\tau$ (s) & coh.\ 0$^\circ$ & coh.\ 90$^\circ$ & $T_2^e$ 0$^\circ$ & $T_2^e$ 90$^\circ$ \\
\midrule
290 & $\ensuremath{2.3\times10^{7}}$ & $\ensuremath{5.7\times10^{-11}}$ & $\ensuremath{1.1\times10^{-14}}$ & $0$ & $0$ & $\ensuremath{-1.0\times10^{-13}}$ \\
160 & $\ensuremath{7.0\times10^{6}}$ & $\ensuremath{5.7\times10^{-11}}$ & $0$ & $\ensuremath{2.2\times10^{-14}}$ & $0$ & $\ensuremath{-8.9\times10^{-14}}$ \\
100 & $\ensuremath{2.7\times10^{6}}$ & $\ensuremath{5.9\times10^{-11}}$ & $0$ & $0$ & $0$ & $\ensuremath{-1.3\times10^{-13}}$ \\
30 & $\ensuremath{2.4\times10^{5}}$ & $\ensuremath{9.3\times10^{-11}}$ & $\ensuremath{-2.2\times10^{-14}}$ & $\ensuremath{-2.2\times10^{-14}}$ & $0$ & $\ensuremath{-5.6\times10^{-13}}$ \\
22 & $\ensuremath{1.3\times10^{5}}$ & $\ensuremath{1.3\times10^{-10}}$ & $\ensuremath{-6.7\times10^{-14}}$ & $\ensuremath{-4.4\times10^{-14}}$ & $\ensuremath{-2.2\times10^{-14}}$ & $\ensuremath{-1.3\times10^{-12}}$ \\
10 & $\ensuremath{2.7\times10^{4}}$ & $\ensuremath{3.9\times10^{-10}}$ & $\ensuremath{-1.6\times10^{-11}}$ & $\ensuremath{-1.8\times10^{-11}}$ & $\ensuremath{-1.6\times10^{-11}}$ & $\ensuremath{-3.5\times10^{-11}}$ \\
3 & $\ensuremath{2.4\times10^{3}}$ & $\ensuremath{3.8\times10^{-9}}$ & $\ensuremath{-8.7\times10^{-8}}$ & $\ensuremath{-7.4\times10^{-7}}$ & $\ensuremath{-8.7\times10^{-8}}$ & $\ensuremath{-7.3\times10^{-7}}$ \\
1 & $\ensuremath{2.7\times10^{2}}$ & $\ensuremath{3.4\times10^{-8}}$ & $\ensuremath{-1.5\times10^{-5}}$ & $\ensuremath{3.9\times10^{-6}}$ & $\ensuremath{-1.5\times10^{-5}}$ & $\ensuremath{3.6\times10^{-6}}$ \\
0.3 & $\ensuremath{2.4\times10^{1}}$ & $\ensuremath{3.7\times10^{-7}}$ & $\ensuremath{-4.3\times10^{-5}}$ & $\ensuremath{6.0\times10^{-6}}$ & $\ensuremath{-4.7\times10^{-5}}$ & $\ensuremath{-2.8\times10^{-6}}$ \\
0.1 & $\ensuremath{2.7\times10^{0}}$ & $\ensuremath{3.4\times10^{-6}}$ & $\ensuremath{-4.9\times10^{-5}}$ & $\ensuremath{6.8\times10^{-6}}$ & $\ensuremath{-7.2\times10^{-5}}$ & $\ensuremath{-1.3\times10^{-4}}$ \\
0.03 & $\ensuremath{2.4\times10^{-1}}$ & $\ensuremath{3.7\times10^{-5}}$ & $\ensuremath{-5.0\times10^{-5}}$ & $\ensuremath{6.9\times10^{-6}}$ & $\ensuremath{-6.9\times10^{-5}}$ & $\ensuremath{-1.5\times10^{-3}}$ \\
0.01 & $\ensuremath{2.7\times10^{-2}}$ & $\ensuremath{3.4\times10^{-4}}$ & $\ensuremath{-5.0\times10^{-5}}$ & $\ensuremath{6.9\times10^{-6}}$ & $\ensuremath{-2.3\times10^{-5}}$ & $\ensuremath{-9.4\times10^{-3}}$ \\
\bottomrule
\end{tabular}

\end{table}

\begin{table}[h]\centering\footnotesize
\caption{Activationless, $\Delta = \lambda = 1.0$~eV.}\label{tab:locact}
\begin{tabular}{rrrrrrr}
\toprule
$t$ (meV) & $|J|$ (MHz) & $\tau$ (s) & coh.\ 0$^\circ$ & coh.\ 90$^\circ$ & $T_2^e$ 0$^\circ$ & $T_2^e$ 90$^\circ$ \\
\midrule
290 & $\ensuremath{4.1\times10^{7}}$ & $\ensuremath{1.0\times10^{-13}}$ & $0$ & $0$ & $0$ & $\ensuremath{-1.1\times10^{-14}}$ \\
160 & $\ensuremath{1.2\times10^{7}}$ & $\ensuremath{1.0\times10^{-13}}$ & $0$ & $0$ & $\ensuremath{1.1\times10^{-14}}$ & $0$ \\
100 & $\ensuremath{4.8\times10^{6}}$ & $\ensuremath{1.1\times10^{-13}}$ & $0$ & $0$ & $0$ & $0$ \\
30 & $\ensuremath{4.4\times10^{5}}$ & $\ensuremath{1.7\times10^{-13}}$ & $0$ & $0$ & $0$ & $0$ \\
22 & $\ensuremath{2.3\times10^{5}}$ & $\ensuremath{2.3\times10^{-13}}$ & $0$ & $0$ & $\ensuremath{-1.1\times10^{-14}}$ & $\ensuremath{-1.1\times10^{-14}}$ \\
10 & $\ensuremath{4.8\times10^{4}}$ & $\ensuremath{7.1\times10^{-13}}$ & $\ensuremath{2.2\times10^{-14}}$ & $\ensuremath{2.2\times10^{-14}}$ & $0$ & $0$ \\
3 & $\ensuremath{4.4\times10^{3}}$ & $\ensuremath{6.8\times10^{-12}}$ & $\ensuremath{8.3\times10^{-12}}$ & $\ensuremath{9.1\times10^{-12}}$ & $\ensuremath{8.3\times10^{-12}}$ & $\ensuremath{9.1\times10^{-12}}$ \\
1 & $\ensuremath{4.8\times10^{2}}$ & $\ensuremath{6.1\times10^{-11}}$ & $\ensuremath{1.5\times10^{-8}}$ & $\ensuremath{5.1\times10^{-8}}$ & $\ensuremath{1.5\times10^{-8}}$ & $\ensuremath{5.1\times10^{-8}}$ \\
0.3 & $\ensuremath{4.4\times10^{1}}$ & $\ensuremath{6.7\times10^{-10}}$ & $\ensuremath{-2.4\times10^{-5}}$ & $\ensuremath{5.1\times10^{-6}}$ & $\ensuremath{-2.4\times10^{-5}}$ & $\ensuremath{5.1\times10^{-6}}$ \\
0.1 & $\ensuremath{4.8\times10^{0}}$ & $\ensuremath{6.1\times10^{-9}}$ & $\ensuremath{-4.3\times10^{-5}}$ & $\ensuremath{6.9\times10^{-6}}$ & $\ensuremath{-4.3\times10^{-5}}$ & $\ensuremath{6.9\times10^{-6}}$ \\
0.03 & $\ensuremath{4.4\times10^{-1}}$ & $\ensuremath{6.7\times10^{-8}}$ & $\ensuremath{-5.2\times10^{-5}}$ & $\ensuremath{6.9\times10^{-6}}$ & $\ensuremath{-5.2\times10^{-5}}$ & $\ensuremath{6.4\times10^{-6}}$ \\
0.01 & $\ensuremath{4.8\times10^{-2}}$ & $\ensuremath{6.1\times10^{-7}}$ & $\ensuremath{-5.0\times10^{-5}}$ & $\ensuremath{6.6\times10^{-6}}$ & $\ensuremath{-5.4\times10^{-5}}$ & $\ensuremath{-1.2\times10^{-5}}$ \\
\bottomrule
\end{tabular}

\end{table}

\section{Partner hyperfine coupling and field orientation}
Coherent MFE (\%) of the MAO contact pair with $J = 0$ and long lifetimes, as a
function of the largest partner hyperfine coupling, with the point-dipole $D$.
This is the most favourable case the contact dipolar coupling allows; it ignores
the competence bound. The same bound with the spin-distributed $D$ is in
Table~\ref{tab:partnerspin}.

\begin{table}[h]\centering\small
\caption{Dependence on the partner hyperfine coupling.}\label{tab:partner}
\begin{tabular}{rrrrr}
\toprule
partner $a$ (MHz) & $\tau$ ($\mu$s) & 0$^\circ$ & 45$^\circ$ & 90$^\circ$ \\
\midrule
5 & 1 & $\ensuremath{-2.4\times10^{-7}}$ & $\ensuremath{-1.2\times10^{-7}}$ & $\ensuremath{3.6\times10^{-8}}$ \\
5 & 100 & $\ensuremath{-2.8\times10^{-7}}$ & $\ensuremath{-1.2\times10^{-7}}$ & $\ensuremath{3.6\times10^{-8}}$ \\
44 & 1 & $\ensuremath{1.5\times10^{-6}}$ & $\ensuremath{9.4\times10^{-7}}$ & $\ensuremath{2.9\times10^{-7}}$ \\
44 & 100 & $\ensuremath{1.5\times10^{-6}}$ & $\ensuremath{9.1\times10^{-7}}$ & $\ensuremath{2.8\times10^{-7}}$ \\
100 & 1 & $\ensuremath{1.5\times10^{-6}}$ & $\ensuremath{4.8\times10^{-6}}$ & $\ensuremath{1.3\times10^{-6}}$ \\
100 & 100 & $\ensuremath{1.6\times10^{-6}}$ & $\ensuremath{4.8\times10^{-6}}$ & $\ensuremath{1.3\times10^{-6}}$ \\
257 & 1 & $\ensuremath{-5.0\times10^{-5}}$ & $\ensuremath{-1.8\times10^{-5}}$ & $\ensuremath{6.7\times10^{-6}}$ \\
257 & 100 & $\ensuremath{-5.0\times10^{-5}}$ & $\ensuremath{-1.8\times10^{-5}}$ & $\ensuremath{6.9\times10^{-6}}$ \\
363 & 1 & $\ensuremath{-8.8\times10^{-5}}$ & $\ensuremath{-3.7\times10^{-5}}$ & $\ensuremath{1.1\times10^{-5}}$ \\
363 & 100 & $\ensuremath{-8.8\times10^{-5}}$ & $\ensuremath{-3.7\times10^{-5}}$ & $\ensuremath{1.1\times10^{-5}}$ \\
\bottomrule
\end{tabular}

\end{table}


\section{Dipolar coupling from the spin distribution}
At van der Waals contact the point-dipole approximation concentrates each
radical's spin on one point. We summed the dipolar tensor over the L\"owdin atomic
spin populations of the two radicals (UKS B3LYP/def2-SVP, vertical ions at the
contact geometry): $T_{ab} = d\sum_{ij}\rho_i\rho_j(\delta_{ab}-3\hat r_a\hat
r_b)/r_{ij}^3$, $d = 52.04$~MHz\,nm$^3$, and took $D = \tfrac34T_n$ from the
principal value of largest magnitude, which reproduces $D = -77.9/r^3$~MHz for a
single pair of unit spins (Table~\ref{tab:dipolar}). The resulting bound, with
$J = 0$ and long lifetimes, is in Table~\ref{tab:partnerspin}.

\begin{table}[h]\centering\small
\caption{Point-dipole and spin-distributed dipolar couplings at the contact
geometry.}\label{tab:dipolar}
\begin{tabular}{rrrrrr}
\toprule
$r$ (\AA) & $D_{\rm point}$ (MHz) & centroid dist.\ (nm) & $D_{\rm spin}$ (MHz) & $E$ (MHz) & principal values (MHz) \\
\midrule
3.50 & -1817 & 0.348 & -914 & -26 & -1219, 558, 662 \\
4.00 & -1217 & 0.397 & -674 & -17 & -899, 415, 484 \\
\bottomrule
\end{tabular}

\end{table}

\begin{table}[h]\centering\footnotesize
\caption{Coherent MFE (\%) with the spin-distributed $D$, $J = 0$.}\label{tab:partnerspin}
\begin{tabular}{lrrrrr}
\toprule
enzyme & partner $a$ (MHz) & $\tau$ ($\mu$s) & 0$^\circ$ & 45$^\circ$ & 90$^\circ$ \\
\midrule
MAO & 5 & 1 & $\ensuremath{-3.4\times10^{-6}}$ & $\ensuremath{-1.7\times10^{-6}}$ & $\ensuremath{5.6\times10^{-7}}$ \\
MAO & 5 & 10 & $\ensuremath{-3.9\times10^{-6}}$ & $\ensuremath{-1.8\times10^{-6}}$ & $\ensuremath{5.6\times10^{-7}}$ \\
MAO & 5 & 100 & $\ensuremath{-4.0\times10^{-6}}$ & $\ensuremath{-1.8\times10^{-6}}$ & $\ensuremath{5.6\times10^{-7}}$ \\
MAO & 44 & 1 & $\ensuremath{6.2\times10^{-6}}$ & $\ensuremath{1.5\times10^{-5}}$ & $\ensuremath{4.3\times10^{-6}}$ \\
MAO & 44 & 10 & $\ensuremath{6.4\times10^{-6}}$ & $\ensuremath{1.6\times10^{-5}}$ & $\ensuremath{4.2\times10^{-6}}$ \\
MAO & 44 & 100 & $\ensuremath{6.4\times10^{-6}}$ & $\ensuremath{1.6\times10^{-5}}$ & $\ensuremath{4.2\times10^{-6}}$ \\
MAO & 100 & 1 & $\ensuremath{-1.0\times10^{-4}}$ & $\ensuremath{-2.5\times10^{-5}}$ & $\ensuremath{1.7\times10^{-5}}$ \\
MAO & 100 & 10 & $\ensuremath{-1.0\times10^{-4}}$ & $\ensuremath{-2.5\times10^{-5}}$ & $\ensuremath{1.8\times10^{-5}}$ \\
MAO & 100 & 100 & $\ensuremath{-1.0\times10^{-4}}$ & $\ensuremath{-2.5\times10^{-5}}$ & $\ensuremath{1.8\times10^{-5}}$ \\
MAO & 257 & 1 & $\ensuremath{-4.3\times10^{-4}}$ & $\ensuremath{-1.9\times10^{-4}}$ & $\ensuremath{6.4\times10^{-5}}$ \\
MAO & 257 & 10 & $\ensuremath{-4.3\times10^{-4}}$ & $\ensuremath{-1.9\times10^{-4}}$ & $\ensuremath{6.6\times10^{-5}}$ \\
MAO & 257 & 100 & $\ensuremath{-4.3\times10^{-4}}$ & $\ensuremath{-1.9\times10^{-4}}$ & $\ensuremath{6.6\times10^{-5}}$ \\
MAO & 363 & 1 & $\ensuremath{-3.3\times10^{-4}}$ & $\ensuremath{-1.3\times10^{-4}}$ & $\ensuremath{7.7\times10^{-5}}$ \\
MAO & 363 & 10 & $\ensuremath{-3.3\times10^{-4}}$ & $\ensuremath{-1.3\times10^{-4}}$ & $\ensuremath{7.9\times10^{-5}}$ \\
MAO & 363 & 100 & $\ensuremath{-3.3\times10^{-4}}$ & $\ensuremath{-1.3\times10^{-4}}$ & $\ensuremath{7.9\times10^{-5}}$ \\
DAO & 5 & 1 & $\ensuremath{-1.1\times10^{-5}}$ & $\ensuremath{-5.4\times10^{-6}}$ & $\ensuremath{1.9\times10^{-6}}$ \\
DAO & 5 & 10 & $\ensuremath{-1.3\times10^{-5}}$ & $\ensuremath{-5.7\times10^{-6}}$ & $\ensuremath{1.9\times10^{-6}}$ \\
DAO & 5 & 100 & $\ensuremath{-1.3\times10^{-5}}$ & $\ensuremath{-5.7\times10^{-6}}$ & $\ensuremath{1.9\times10^{-6}}$ \\
DAO & 58 & 1 & $\ensuremath{-8.6\times10^{-5}}$ & $\ensuremath{3.5\times10^{-6}}$ & $\ensuremath{2.2\times10^{-5}}$ \\
DAO & 58 & 10 & $\ensuremath{-8.6\times10^{-5}}$ & $\ensuremath{4.3\times10^{-6}}$ & $\ensuremath{2.2\times10^{-5}}$ \\
DAO & 58 & 100 & $\ensuremath{-8.6\times10^{-5}}$ & $\ensuremath{4.3\times10^{-6}}$ & $\ensuremath{2.2\times10^{-5}}$ \\
DAO & 100 & 1 & $\ensuremath{-4.1\times10^{-4}}$ & $\ensuremath{-1.5\times10^{-4}}$ & $\ensuremath{5.3\times10^{-5}}$ \\
DAO & 100 & 10 & $\ensuremath{-4.1\times10^{-4}}$ & $\ensuremath{-1.5\times10^{-4}}$ & $\ensuremath{5.5\times10^{-5}}$ \\
DAO & 100 & 100 & $\ensuremath{-4.1\times10^{-4}}$ & $\ensuremath{-1.5\times10^{-4}}$ & $\ensuremath{5.5\times10^{-5}}$ \\
\bottomrule
\end{tabular}

\end{table}

\section{Cryptochrome control}
FAD$^{\bullet-}$--W324$^{\bullet+}$ at 1.90~nm, $D = \DCRY$~MHz, lifetime and
$T_2^e$ of 1~$\mu$s, exchange from $J(r) = J_0 e^{-\beta r}$ with
$\beta = 14$~nm$^{-1}$ over $8\times10^{12}\le|J_0|/\mu{\rm T}\le8\times10^{14}$ and
both signs (Efimova and Hore). Spin-$\tfrac12$: N5 and N10 as effective
spin-$\tfrac12$ nuclei; spin-1: $^{14}$N as $I = 1$. Trp H$\beta$1 is 45~MHz in
both. Powder: Gauss--Legendre average over the field direction.

\begin{table}[h]\centering\small
\caption{Cryptochrome MFE (\%) at 50~$\mu$T.}\label{tab:cry}
\begin{tabular}{llrrrr}
\toprule
model & $|J_0|$ ($\mu$T) & $J$ (MHz) & 0$^\circ$ & 90$^\circ$ & powder \\
\midrule
spin half & $\ensuremath{8\times10^{12}}$ & -1.26 & -0.236 & -1.742 & -1.422 \\
spin one N & $\ensuremath{8\times10^{12}}$ & -1.26 & +0.186 & -1.447 & -- \\
spin half & $\ensuremath{8\times10^{12}}$ & +1.26 & +0.199 & -1.048 & -0.880 \\
spin one N & $\ensuremath{8\times10^{12}}$ & +1.26 & +0.262 & -1.285 & -- \\
spin half & $\ensuremath{8\times10^{13}}$ & -12.57 & -0.025 & -0.720 & -0.501 \\
spin one N & $\ensuremath{8\times10^{13}}$ & -12.57 & -0.016 & -1.436 & -- \\
spin half & $\ensuremath{8\times10^{13}}$ & +12.57 & +0.058 & -0.284 & -0.177 \\
spin one N & $\ensuremath{8\times10^{13}}$ & +12.57 & +0.029 & -0.206 & -- \\
spin half & $\ensuremath{8\times10^{14}}$ & -125.73 & -0.012 & -0.322 & -0.254 \\
spin one N & $\ensuremath{8\times10^{14}}$ & -125.73 & -0.037 & -0.222 & -- \\
spin half & $\ensuremath{8\times10^{14}}$ & +125.73 & -0.021 & -0.326 & -0.222 \\
spin one N & $\ensuremath{8\times10^{14}}$ & +125.73 & -0.022 & -0.234 & -- \\
\bottomrule
\end{tabular}

\end{table}

\end{document}
